% !PRINTED 687
\weberpage{687}
By (9), $\alpha_1$, $\beta_1$ vanish at the same points and with the same order numbers as $\alpha$, $\beta$, and from this there follow, corresponding to the representations (2) and (4), the polygon representations:
\weq{11}{\alpha_1 = \dfrac{\mathfrak{M}}{\mathfrak{A}_1}, \quad \beta_1 = \dfrac{\mathfrak{N}}{\mathfrak{B}_1},}
and from (5), (7) and (10):
\weq{12}{F_1'(\alpha_1) = \dfrac{\mathfrak{L}}{\mathfrak{A}_1^{n-2}\mathfrak{B}_1^m}.}

With this, formula (8) is proved \gesperrt{in general}.

If one sets, following \S\ 194, (9):
\[
F'(\alpha)\,d\alpha + F'(\beta)\,d\beta = 0,
\]
then from (8) and \S\ 195, (12) it follows that:
\weq{13}{F'(\beta) = \dfrac{\mathfrak{R}\mathfrak{Z}_\alpha}{\mathfrak{A}^n\mathfrak{B}^{m-2}}.}

Both derivatives $F'(\alpha)$, $F'(\beta)$ therefore vanish at the points of $\mathfrak{R}$, and $\mathfrak{R}$ is called the polygon of the double points in $(\alpha,\beta)$.

The order of $\mathfrak{A}^n\mathfrak{B}^{m-2}$ is $2n(m-1)$, and from this there results for the order $2r$ of $\mathfrak{R}$ an even number, namely:
\[
2r = 2n(m-1) - w_\alpha = 2m(n-1) - w_\beta,
\]
and for the genus $p$ there results, by (2):
\weq{14}{p = (n-1)(m-1) - r.}

\begin{center}
\rule{4cm}{0.4pt}
\end{center}
